% Felipe Muñoz Casasola
% This is free and unencumbered software released into the public domain.

% Anyone is free to copy, modify, publish, use, compile, sell, or
% distribute this software, either in source code form or as a compiled
% binary, for any purpose, commercial or non-commercial, and by any
% means.

% In jurisdictions that recognize copyright laws, the author or authors
% of this software dedicate any and all copyright interest in the
% software to the public domain. We make this dedication for the benefit
% of the public at large and to the detriment of our heirs and
% successors. We intend this dedication to be an overt act of
% relinquishment in perpetuity of all present and future rights to this
% software under copyright law.

% THE SOFTWARE IS PROVIDED "AS IS", WITHOUT WARRANTY OF ANY KIND,
% EXPRESS OR IMPLIED, INCLUDING BUT NOT LIMITED TO THE WARRANTIES OF
% MERCHANTABILITY, FITNESS FOR A PARTICULAR PURPOSE AND NONINFRINGEMENT.
% IN NO EVENT SHALL THE AUTHORS BE LIABLE FOR ANY CLAIM, DAMAGES OR
% OTHER LIABILITY, WHETHER IN AN ACTION OF CONTRACT, TORT OR OTHERWISE,
% ARISING FROM, OUT OF OR IN CONNECTION WITH THE SOFTWARE OR THE USE OR
% OTHER DEALINGS IN THE SOFTWARE.

% For more information, please refer to <https://unlicense.org/>

\documentclass{article}

\usepackage[spanish]{babel}
\usepackage[margin=3cm]{geometry}
\usepackage{microtype}
\usepackage[pdfusetitle]{hyperref}
\usepackage{minted}
\usepackage{booktabs}

\title{Gestor de archivos lf}
\author{Felipe Muñoz Casasola}
\date{22 de febrero de 2026}

\begin{document}
    \maketitle
    \clearpage

    \section{Instalación}
    Lf se instala desde los repositorios, aunque esté desactualizado.

    \section{Configuración}
    Los archivos de configuración se encuentran en \texttt{$\sim$/.config/lf},
    actualmente son los archivos \texttt{lfrc}, \texttt{icons} y \texttt{colors}.
    El archivo \texttt{colors} es el que viene por defecto del repositorio de
    \texttt{lf}. En el que hay unos archivos de configuración de ejemplo. Es
    recomendable revisarlos alguna vez al reinstalarlos, así como la
    documentación de \texttt{lf}. A los que se puede acceder a través de su
    repositorio oficial.

    \section{Uso}

    \begin{table}[h]
    	   \begin{tabular}{cc}
    	   	   \midrule
		   Atajo	&	Acción	\\
		   \midrule

		   \texttt{→}	&	Abrir archivo o directorio actual.	\\
           \texttt{←}   &   Subir al directorio padre.  \\
           \texttt{q}    &   Salir. \\
		   \texttt{Espacio}	&	Seleccionar archivo o directorio actual.	\\
		   \texttt{Ct}	&	Comprimir en un tar (dentro del tar se crea una carpeta comprimidos).	\\
		   \texttt{Cz}	&	Comprimir en un zip (dentro del zip se crea una carpeta comprimidos).	\\
           \texttt{E}	&	Extraer en el directorio actual.	\\
		   \texttt{Co}	&	Cortar la selección.	\\
		   \texttt{y}	&	Copiar la selección.	\\
		   \texttt{r}	&	Renombrar.	\\
		   \texttt{Supr}	&	Mover a la papelera (puede fallar con directorios).	\\
		   \texttt{id}	&	Ir a $\sim$/Documentos.	\\
           \texttt{ii}	&	Ir a $\sim$/Imágenes.	\\
           \texttt{iu}	&	Ir a los USB (con LXQt ya se me montan por defecto).	\\
           \texttt{ic}	&	Ir a $\sim$/Escritorio/Compartir.	\\
           \texttt{it}	&	Ir a /tmp.	\\
           \texttt{im}	&	Ir a $\sim$/Descargas.	\\
           \texttt{m}   &   Mostrar archivos ocultos.	\\
		   \midrule
    	   \end{tabular}
    \end{table}

    \section{Archivos de configuración}
    Archivo \texttt{lfrc}:

    \inputminted[breaklines, breakanywhere, breaksymbolleft="", breaksymbolright="", breakanywheresymbolpre="", breakanywheresymbolpost=""]
    {vim}{inc/lfrc}

    Archivo \texttt{icons}:

    \inputminted[breaklines, breakanywhere, breaksymbolleft="", breaksymbolright="", breakanywheresymbolpre="", breakanywheresymbolpost=""]
    {vim}{inc/icons}

    Para cambiar los iconos puedes leer el manual del gestor de ventanas, pues
    allí viene cómo insertar iconos en Vim para poder cambiar los iconos.

    Archivo \texttt{colors}:

    \inputminted[breaklines, breakanywhere, breaksymbolleft="", breaksymbolright="", breakanywheresymbolpre="", breakanywheresymbolpost=""]
    {vim}{inc/colors}

    \section{Arreglar errores}
    \subsection{Vim sin resaltado de sintaxis}
    Lo que ocurre es que por defecto \texttt{lf} usa por defecto \texttt{vi}
    como editor de texto. Actualmente se le puede decir que use otro declarando
    la variable \texttt{EDITOR} a \texttt{vim}.

    \subsection{Los iconos no funcionan}
    Asegúrate de haber activado los iconos en el archivo de configuración
    \texttt{lfrc}. Y de tener en la terminal activa la fuente de esos
    iconos, que actualmente es una NerdFont.

    \subsection{Ejecutar \texttt{lf} como superusuario}
    Para lograr esto basta invocar una terminal que ejecute el comando
    \texttt{lf} precedido por \texttt{sudo}.
\end{document}
